\documentclass[12pt,a4paper]{article}

\usepackage[a4paper,margin=1in]{geometry}
\usepackage[T1]{fontenc}
\usepackage[utf8]{inputenc}
\usepackage{enumitem}
\usepackage{tabularx}
\usepackage{array}
\usepackage{longtable}
\usepackage{hyperref}

\setlength{\emergencystretch}{3em}
\newcolumntype{Y}{>{\raggedright\arraybackslash}X}

\title{\textbf{Functional Requirements}\\Tower Defence Game}
\author{Jonas Leipner}
\date{\today}

\begin{document}

\maketitle

\section{Introduction}
\label{sec:introduction}

This document translates the project's initial user stories into functional
requirements. The requirements are organised around three personas:
\begin{itemize}[leftmargin=*]
    \item \textbf{Player:} the person who plays the tower defence game,
    places towers, and responds to incoming enemies.
    \item \textbf{Researcher:} the person who analyses player behaviour and
    uses telemetry to investigate player strategies.
    \item \textbf{System Architecture:} the architectural responsibilities
    concerned with extracting game state and executing game logic dynamically.
\end{itemize}

The project is being developed iteratively using an Agile/Scrum-inspired
approach. The requirements should therefore be treated as a prioritised
product backlog: Must Have items define the initial product increment, while
Should Have and Could Have items can be scheduled into later iterations if
they remain compatible with the project's time and research scope.

\section{MoSCoW Prioritisation}
\label{sec:moscow}

The MoSCoW method is used to control scope:
\begin{description}[leftmargin=2.7cm,style=nextline]
    \item[Must Have] Essential for the minimum viable product or for testing
    the central architectural experiment.
    \item[Should Have] Important and valuable, but the system can still
    function without it in the first product increment.
    \item[Could Have] Desirable enhancements that are implemented only if
    sufficient time and resources remain.
    \item[Won't Have] Explicitly excluded from the current project scope,
    although it may be reconsidered in future work.
\end{description}

\section{Must Have Requirements}
\label{sec:must-have}

\subsection{Player Persona}

\begin{description}[leftmargin=1.2cm,style=nextline]
    \item[FR-M1 --- Tower placement:] The system shall allow the player to
    select towers and place a maximum of five towers on the game grid. The
    placed towers shall be usable to defend the home base from incoming
    monsters.
\end{description}

This requirement defines the basic game loop. Without tower placement and a
limit on the number of towers, there is no meaningful defensive gameplay.

\subsection{System Architecture Persona}

\begin{description}[leftmargin=1.2cm,style=nextline]
    \item[FR-M2 --- Dynamic combat logic:] The system shall extract the
    current combat state from the Flame game engine and pass that state to a
    \texttt{dart\_eval} sandbox. The sandbox shall be capable of executing
    game logic dynamically without requiring an application restart or a new
    ahead-of-time compilation.
\end{description}

This requirement is the central architectural experiment of the project. The
game may be playable without dynamic evaluation, but the project objective
depends on demonstrating this Flame-to-\texttt{dart\_eval} boundary.

\section{Should Have Requirements}
\label{sec:should-have}

\subsection{Player Persona}

\begin{description}[leftmargin=1.2cm,style=nextline]
    \item[FR-S1 --- Enemy scoring:] The system should assign different point
    values to different enemy types. The player should be able to use the
    resulting score to measure success as game waves become progressively more
    difficult.
\end{description}

Different enemy values improve feedback, enjoyment, and progression. However,
the core defensive game loop can function without a fully developed scoring
system.

\subsection{Researcher Persona}

\begin{description}[leftmargin=1.2cm,style=nextline]
    \item[FR-S2 --- Tower telemetry:] The system should capture telemetry data
    describing which towers are placed most frequently. The data should support
    analysis of player strategies and provide a possible foundation for a
    future dynamic recommender system.
\end{description}

Telemetry supports the research goals but is not required for the game to be
playable. The first implementation should therefore remain lightweight, such
as a local structure recording tower-placement events.

\section{Could Have Requirements}
\label{sec:could-have}

\subsection{Player Persona}

\begin{description}[leftmargin=1.2cm,style=nextline]
    \item[FR-C1 --- Tower merging:] The system could allow the player to merge
    two towers in an incubator to create a unique tower. This feature could
    allow players to discover new defensive options and achieve a higher score.
\end{description}

Tower merging could provide a distinctive gameplay mechanic, but it introduces
additional algorithmic, interface, balancing, and testing complexity. It is
therefore outside the immediate minimum viable product.

\subsection{System Architecture Persona}

\begin{description}[leftmargin=1.2cm,style=nextline]
    \item[FR-C2 --- Local outcome logging:] The system could log player
    choices and victory metrics locally. The logged data could later be
    formatted and supplied to a dynamic recommender system.
\end{description}

This requirement is a possible future research direction rather than a fixed
part of the initial project scope. It can be considered after the core
architecture and gameplay loop have been demonstrated.

\section{Won't Have Requirements}
\label{sec:wont-have}

No requirements have currently been assigned to the Won't Have category.
Requirements placed in this category would be explicitly excluded from the
current project scope and would not be scheduled into the planned increments.

\section{Prioritised Product Backlog}
\label{sec:backlog}

Table~\ref{tab:requirements} provides a compact backlog view for planning
iterations or Scrum sprints.

\begin{longtable}{>{\raggedright\arraybackslash}p{0.12\textwidth}
                  >{\raggedright\arraybackslash}p{0.17\textwidth}
                  >{\raggedright\arraybackslash}p{0.63\textwidth}}
    \caption{Prioritised functional requirements}
    \label{tab:requirements}\\
    \hline
    \textbf{ID} & \textbf{Priority} & \textbf{Requirement} \\
    \hline
    \endfirsthead
    \hline
    \textbf{ID} & \textbf{Priority} & \textbf{Requirement} \\
    \hline
    \endhead
    FR-M1 & Must Have & Place a maximum of five towers on the grid to defend
    the home base. \\
    FR-M2 & Must Have & Pass combat state from Flame to a
    \texttt{dart\_eval} sandbox for dynamic game logic. \\
    FR-S1 & Should Have & Assign different point values to enemy types and
    use them for player scoring. \\
    FR-S2 & Should Have & Capture telemetry about frequently placed towers. \\
    FR-C1 & Could Have & Merge two towers to create a unique tower. \\
    FR-C2 & Could Have & Log player choices and victory metrics locally. \\
    \hline
\end{longtable}

\section{Agile Use}
\label{sec:agile-use}

At the start of each iteration, the team can select a small number of
requirements from the backlog and break them into implementation tasks, tests,
and acceptance criteria. A requirement should be considered complete only when
its behaviour has been implemented, tested, and demonstrated. The backlog can
be updated after each review if project findings or supervisor feedback change
the priority of a requirement.

\end{document}
